\documentclass[../main.tex]{subfiles}
\graphicspath{{\subfix{../img/}}}
\begin{document}
	\subsection{Campo Magnetico}
	Il fenomeno del magnetismo era noto in grecia fin dall’ 800 a.c.\\
	Magnetite ($F_{e_3}O_4$) [magnesia-grecia] minerale di cui era nota la proprietà di esercitare forze su minerali simili e su pezzi di ferro.\\\\
	Ogni magnete, indipendentemente dalla sua forma, mostra la presenza di due poli, ossia di due parti in cui si localizzano le proprietà di attrazione sono petti \textbf{polo nord} e \textbf{polo sud}.
	\begin{center}
		\begin{figure}[h]
			\centering
			\includegraphics[scale=0.55]{92.png}
			\caption{}
			\label{fig:im-92}
		\end{figure}
	\end{center}
	\begin{itemize}
		\item La forza è attrattiva o repulsiva a seconda chei poli siano opposti o uguali
		\item In un magnete i poli sono sempre di segno \textbf{opposto}
	\end{itemize}
	$$\Downarrow$$
	Non è possibile ottenere un polo magnetico isolato $\rightarrow$ esistono solo come dipoli magnetici.
	Esempio: se spezziamo un magnete:
	\begin{center}
		\begin{figure}[h]
			\centering
			\includegraphics[scale=0.55]{93.png}
			\caption{Magnete spezzato}
			\label{fig:im-93}
		\end{figure}
	\end{center}
	\begin{itemize}
		\item Oggetti in ferro oltre ad essere attratti dai magneti possono diventare essi stessi dei magneti artificiali o calamite
		\item Un magnete in forma di sbarretta sottile (ago magnetico) sospeso e libero di ruotare in un piano orizzontale tende ad allinearsi parallelamente al meridiano terrestre
	\end{itemize} 
	\begin{center}
		\begin{figure}[h]
			\centering
			\includegraphics[scale=0.45]{94.png}
			\caption{}
			\label{fig:im-94}
		\end{figure}
	\end{center}
	Dal punto di vista dei dati sperimentali possiamo anche anticipare che nel 1811 fu scoperta una relazione fra fenomeni magnetici ed elettrici
	\begin{center}
		\begin{figure}[h]
			\centering
			\includegraphics[scale=0.45]{95.png}
			\caption{}
			\label{fig:im-95}
		\end{figure}
	\end{center}
	Un ago magnetico posto nelle vicinanze di un filo percorso da corrente assume una ben definita posizione di equilibrio
	$\Rightarrow$ cariche elettriche in movimento come sorgenti di azioni magnetiche.\\\\
	In elettrostatica per descrivere le “azioni elettriche” fra cariche stazionarie o in moto si era introdotto il concetto di campo elettrico nei punti dello spazio.\\
	Ora dobbiamo introducere un “analogo” magnetico per poter descrivere le azioni magnetiche
	\subsubsection{Il Campo Magnetico $\vec{B}$}
	Per descrivere lo spazio circostante una sostanza magnetizzata introduciamo un vettore $\vec{B}$ per rappresentare intensità e orientazione del campo magnetico.\\
	Analogamente a quanto visto per il campo elettrico rappresentiamo il campo magnetico per via grafica attraverso le linee del campo magnetico:
	\begin{center}
		\begin{figure}[h]
			\centering
			\includegraphics[scale=0.45]{96.png}
			\caption{Linee campo magnetico}
			\label{fig:im-96}
		\end{figure}
	\end{center}
	\begin{itemize}
		\item In ogni suo punto una linea di forza è tangente e concorde al campo $\vec{B}$ in quel punto
		\item Si addensano dove l’intensità del campo è maggiore
		\item Non si incrociano mai
		\item Hanno origine dal polo nord e terminano in quello sud
	\end{itemize}
	\textbf{L’orientazione} del campo magnetico in un punto dello spazio coincide con l’orientazione di un ago magnetico posto in quel punto.\\\\
	\subsubsection{Come definiamo $\vec{B}$? (Forza di Lorentz)}
	Come per il campo elettrico si parte dalla misura di una forza agente su un opportuno “oggetto di prova”.\\
	$\rightarrow F_b$: Forza magnetica\\
	$\rightarrow$ “oggetto di prova”: particella di carica $q$ e velocità $\vec{v}$\\
	la forza magnetica agente su una particella carica di carica $q$ e velocità $\vec{v}$ si chiama \textbf{Forza di Lorentz}.
	\begin{center}
		\begin{figure}[h]
			\centering
			\includegraphics[scale=0.45]{97.png}
			\caption{}
			\label{fig:im-97}
		\end{figure}
	\end{center}
	$$\vec{F}_B = q \vec{v} \times \vec{B}$$
	$\Rightarrow$ in modulo $F_B = q \; v \; B \; \sin \vartheta$
	$\Rightarrow$ nel S.I. $[B] \rightarrow$ Tesla (T)
	$$1 T = 1 \frac{N}{C \cdot m/s} = 1 \frac{N}{A \cdot m}$$
	$$1 T = 10^4 G \Leftrightarrow 1 \text{ Tesla} = 10^4 \text{ Gauss}$$
	$$\vec{F}_B = q \vec{v} \times \vec{B}$$
	\begin{figure}[h]
		
		\begin{subfigure}{0.33\textwidth}
			\includegraphics[scale=0.55]{98.png} 
			\caption{}
			\label{fig:im-98}
		\end{subfigure}
		\begin{subfigure}{0.33\textwidth}
			\includegraphics[scale=0.55]{99.png}
			\caption{}
			\label{fig:im-99}
		\end{subfigure}
		\begin{subfigure}{0.33\textwidth}
			\includegraphics[scale=0.55]{100.png}
			\caption{Verso opposto se $q < 0$}
			\label{fig:im-100}
		\end{subfigure}
		
		\caption{}
		\label{fig:image8}
	\end{figure}
	\begin{itemize}
		\item In presenza di un campo magnetico $\vec{B}$ la forza magnetica agente su una particella è proporzionale alla carica $q$ e alla velocità $\vec{v}$
		\item È un prodotto vettoriale!!!
		\item La forza $\vec{f}_B$ si inverte se cambiamo il segno della carica
		\item • Se $\vec{B} \parallel \vec{v}$ la forza magnetica è nulla $F_B = q v B \sin \theta$
	\end{itemize}
	\begin{center}
		\begin{tabular}{|l|l|}
			\hline
			\textbf{Forza Elettrica} & \textbf{Forza Magnetica} \\ 
			\hline
			$\vec{F}_E = q \vec{E}$ & $\vec{F}_B = q \vec{v} \times \vec{B}$\\
			\hline
			$\bullet$ è sempre $\parallel$ ad $\vec{E}$ & $\bullet$  è sempre $\perp$ a $\vec{B}$ \\
			$\bullet$ è \textbf{indipendente} dalla velocità $v$ della particella & $\bullet$ \textbf{dipende} dalla velocità della particella \\
			$\bullet$ compie lavoro quando agisce sulla particella & $\bullet$ \textbf{non} compie lavoro su una particella che si sposta ($\vec{B}$ Statico) \\
			$W = \int_A^B \vec{F} \cdot d\vec{s} = q \int_A^B \vec{E} \cdot d\vec{s} = q (V_A - V_B)$ & $\Downarrow$ \\
			\quad & la forza magnetica $\vec{F}_B$ è infatti sempre $\perp$ alla velocità $\vec{v}$ \\
			\quad & $\Downarrow$ \\
			\quad & $W = \int_A^B \vec{F} \cdot d\vec{s} = q \int_A^B \underbrace{(\vec{v} \times \vec{e})}_{1 \cdot d\vec{s}} \cdot d\vec{s} = 0$ \\
			\hline
		\end{tabular}
	\end{center}
	$\Rightarrow$ la velocità per effetto della forza magnetica cambierà \textbf{solo in direzione e non in modulo}\\
	$\Rightarrow$ l’energia cinetica \textbf{non cambia}\\
	$\Rightarrow$ non abbiamo accelerazione tangenziale ma solo \textbf{centripeta}
	\subsubsection{Esempi}
	\paragraph{Esempio 1}
	\begin{center}
		\begin{figure}[h]
			\centering
			\includegraphics[scale=0.45]{101.png}
			\caption{}
			\label{fig:im-101}
		\end{figure}
	\end{center}
	Una particella carica ($+q$) si muove in un campo \textbf{uniforme} $\vec{B}$ con $\vec{B} \perp \vec{v}$
	$$\Rightarrow \vartheta = \frac{\pi}{2} \Rightarrow \sin \vartheta = 1$$
	Dove $\vartheta$ è \textbf{l'angolo formato con il campo magnetico}, che in questo caso è $\frac{\pi}{2}$, ortogonale ad esso.
	$$F_B = q v B = m \; a_c = m \frac{v^2}{r}$$
	Dove $r = \frac{m v}{q B}$ è il \textbf{Raggio di curvatura}\\\\
	$\Rightarrow$ la traettoria è un arco di circonferenza se $\vec{B}$ uniforme $\Rightarrow r = \text{COSTANTE}$ \\
	Definiamo la \textbf{velocità angolare del moto} come $\omega = \frac{v}{2} = \frac{q B}{m}$ \\
	La frequenza corrispondente alla velocità ancolare del modo è detta \textbf{frequenza di ciclotrone} ed è definita come $\gamma = \frac{\omega}{2\pi} = \frac{q B}{2\pi m}$\\
	La frequenza di rivoluzione non vivende dal raggio e dalla velocità della particella.
	\paragraph{Esempio 2}
	Se $\vartheta$ generico (e $\vec{B}$ uniforme)
	\begin{center}
		\begin{figure}[h]
			\centering
			\includegraphics[scale=0.45]{102.png}
			\caption{}
			\label{fig:im-102}
		\end{figure}
	\end{center}
	$\vec{v} = \vec{v}_\perp + \vec{v}_\parallel \Rightarrow \vec{F}_B = q \vec{v} \times \vec{B} = q (\vec{v}_\perp + \vec{v}_\parallel) \times \vec{B} $
	$= q \vec{v}_\perp \times \vec{B} + q \vec{v}_\parallel \times \vec{B} $
	Il contributo parallelo non interagisce con il campo magnetico, di conseguenza è uguale a zero e possiamo scrivere il tutto come:
	$$= \boxed{q \vec{v}_\perp \times \vec{B}}$$
	\begin{center}
		\begin{figure}[h]
			\centering
			\includegraphics[scale=0.45]{103.png}
			\caption{}
			\label{fig:im-103}
		\end{figure}
	\end{center}
	$$r = \frac{m v_\perp}{q B} \qquad \text{ il moto è elicoidale.}$$ \\
	P: \textbf{passo dell’elica}
	$$P = v_\parallel T = v_\parallel \frac{1}{\gamma} = v_\parallel \frac{2\pi m}{q B} = 2\pi r \frac{v_\parallel}{v_\perp}$$
	\paragraph{Esempio 3}
	Se nello spazio sono presenti contemporaneamete un campo elettrico $\vec{E}$ ed un campo magnetico $\vec{B}$
	$$\Rightarrow \vec{F} = \vec{F}_E + \vec{F}_B = q\vec{E} + q\vec{v} \times \vec{B} \quad \text{somma vettoriale}$$
	Una applicazione tipica si ha con $\vec{E} \perp \vec{B}$ $\qquad$($E\vec{v} \perp \vec{E}, \vec{B}$)
	\begin{figure}[h]
		
		\begin{subfigure}{0.55\textwidth}
			\includegraphics[scale=0.55]{104.png} 
			\caption{}
			\label{fig:im-104}
		\end{subfigure}
		\begin{subfigure}{0.55\textwidth}
			\includegraphics[scale=0.55]{105.png}
			\caption{}
			\label{fig:im-105}
		\end{subfigure}
		
		\caption{}
		\label{fig:image9}
	\end{figure}
	La deflessione della particella $\vec{B}$ nulla quando $\vec{B}$ nulla la forza che agisce su di essa $\Rightarrow qE = qvB \Rightarrow v = \frac{E}{B}$
	In questo modo si può realizzare un selettore di velocità.\\\\
	Se il fascio così ottenuto veniste immesso in una regione di spazio in qui sia presente il solo campo magnetico:
	\begin{center}
		\begin{figure}[h]
			\centering
			\includegraphics[scale=0.45]{106.png}
			\caption{}
			\label{fig:im-106}
		\end{figure}
	\end{center}
	$$r = \frac{mv}{qB_0} \Rightarrow \frac{m}{q} = \frac{r B_0}{ v}$$
	Si realizza uno \textbf{spettrometro di massa}.\\\\
	Ioni isoenergetici di uno stesso elemento seguiranno traettorie diverse in base alla massa $m \Rightarrow$ \textbf{Rivelazione degli isotopi}.
	\subsection{Effetto Hall}
	La forza di Lorentz può deflettere anche i portatori di carica all’interno di un materiale conduttore. \\
	Si consideri una lamina di un materiale conduttore (percorsa da corrente) immersa in un campo magnetico uniforme.
	\begin{figure}[h]
		
		\begin{subfigure}{0.55\textwidth}
			\includegraphics[scale=0.55]{107.png} 
			\caption{}
			\label{fig:im-107}
		\end{subfigure}
		\begin{subfigure}{0.55\textwidth}
			\includegraphics[scale=0.55]{108.png}
			\caption{}
			\label{fig:im-108}
		\end{subfigure}
		
		\caption{}
		\label{fig:image10}
	\end{figure}
	Il campo elettrico che si instaura trasversalmente alla direzione della corrente è detto \textbf{campo di Hall ($E_H$)}.
	La differenza di potenziale è detta \textbf{tensione di Hall}.
	\begin{center}
		\begin{figure}[h]
			\centering
			\includegraphics[scale=0.45]{109.png}
			\caption{}
			\label{fig:im-109}
		\end{figure}
	\end{center}
	Dalla condizione di equilibrio:
	$$F_B = F_{E_H} \Rightarrow \cancel{e} v_d B = \cancel{e} E_H \Rightarrow E_H = v_d B$$
	Dove ovviamente $e$ indica la carica di un elettrone.\\\\
	Ricordando che $J = n \, e \, v_d$ e che $J = \frac{I}{A} = \frac{I}{d \, t}$\\
	Dove $d \, t$ indicano distanza e spessore non differenziale.
	$$\Rightarrow v_d = \frac{J}{n \, e} = \frac{I}{n \, e \, d \, t} \Rightarrow E_H = \frac{I}{n \, e \, d \, t} B \Rightarrow \frac{\Delta V_H}{\cancel{d}} = \frac{I}{n \, e \, \cancel{d} \, t} B$$
	$$\Rightarrow \Delta V_H = \frac{IB}{n \, e \, t} = R_H \frac{IB}{t} $$
	Dove $R_H$ indica il \textbf{coefficiente di Hall}.
	\begin{enumerate}
		\item Possibilità di misurare sperimentalmente $B$ (\textbf{Sonda di Hall})
		\item Il segno di $R_H$ fornisce il segno dei portatori di carica $+$ densità $n$
	\end{enumerate}
	\begin{figure}[h]
		
		\begin{subfigure}{0.55\textwidth}
			\includegraphics[scale=0.55]{110.png} 
			\caption{}
			\label{fig:im-110}
		\end{subfigure}
		\begin{subfigure}{0.55\textwidth}
			\includegraphics[scale=0.55]{111.png}
			\caption{}
			\label{fig:im-111}
		\end{subfigure}
		
		\caption{}
		\label{fig:image11}
	\end{figure}
	\subsection{Forza Magnetica su un Conduttore Percorso da Corrente (Stazionaria)}
	\begin{center}
		\begin{figure}[h]
			\centering
			\includegraphics[scale=0.45]{112.png}
			\caption{}
			\label{fig:im-112}
		\end{figure}
	\end{center}
	In un conduttore percorso da corrente gli elettroni si muovono con velocità $V_d$\\
	$ \Rightarrow \vec{F}_B = -e \vec{v_d} \times \vec{B}$ per ciascun elettrone.\\\\
	Se $n$ è la densità degli elettroni (numero di portatori per unità di volume)\\
	$\Rightarrow$ nel tratto considerato abbiamo $n \, A \, L$ elettroni\\
	$\Rightarrow$ la forza magnetica totale sul tratto di filo è $\vec{F}_B = -e \, \vec{v_d} \times \vec{B} n\, A \, L$\\\\
	Ma in un filo $\vec{J} = -n\, e \, \vec{v_d} \Rightarrow \vec{F}_B = \vec{J} \times \vec{B} \, A \, L$\\
	Con l’introduzione della densità di corrente $J$ si vede che $\vec{F}_B$ non dipende dal segno dei portatori, ma solo dal verso della corrente.\\\\
	Essendo $J = \frac{I}{A}$ per un filo rettilineo possiamo scrivere:
	$$\vec{F}_B = I \vec{L} \times \vec{B}$$
	Dove $\vec{L}$ è un vettore con modulo $L$ e direzione e verso di $\vec{J}$.\\\\
	\subsubsection{Generalizzando per un filo di forma arbitraria}
	\begin{center}
		\begin{figure}[h]
			\centering
			\includegraphics[scale=0.45]{113.png}
			\caption{}
			\label{fig:im-113}
		\end{figure}
	\end{center}
	\begin{itemize}
		\item Percorso da una corrente $I$ 
		\item In un campo magnetico $\vec{B}$
	\end{itemize}
	$$d\vec{F}_B = I \vec{dl} \times \vec{B}$$
	\begin{center}
		\begin{figure}[h]
			\centering
			\includegraphics[scale=0.45]{114.png}
			\caption{}
			\label{fig:im-114}
		\end{figure}
	\end{center}
	La forza su un tratto infinitetimo è ortogonale al campo magnetico e al filo (è uno strumento di calcolo).\\\\
	Per ottenere la forza totale
	$$\vec{F}_B = I \int_a d\vec{l} \times \vec{B}$$
	Che è sempre un prodotto vettoriale!!!
	\subsection{Momento Meccanico su una Spira Piana}
	Consideriamo una spira piana (rettangolare) percorsa da una corrente $I$ e immersa in un campo magnetico uniforme.
	\begin{figure}[h]
		
		\begin{subfigure}{0.55\textwidth}
			\includegraphics[scale=0.55]{115.png} 
			\caption{}
			\label{fig:im-115}
		\end{subfigure}
		\begin{subfigure}{0.55\textwidth}
			\includegraphics[scale=0.55]{116.png}
			\caption{}
			\label{fig:im-116}
		\end{subfigure}
		
		\caption{}
		\label{fig:image12}
	\end{figure}
	Su ciascun lato agisce una forza magnetica $\vec{F}_B = I \int_{a}^{b} d\vec{s} \times \vec{B}$\\
	Che per fili rettilinei in campo uniforme diventa $\vec{F}_B = I \, \vec{L} \times \vec{B}$
	$$F_{AB} = I  \; \overline{AB} \;  B \; \sin(90^\degree + \vartheta) $$
	$$F_{CD} = I  \; \overline{CD} \;  B \; \sin(90^\degree - \vartheta) $$
	Dove possiamo notare che le due forze sono uguali in modulo e giacenti sulla stessa retta di azione.
	$$F_{BC} = I \; \overline{BC} \; B$$
	$$F_{AD} = I \; \overline{AD} \; B$$
	Mentre qui notiamo che sono uguali in modulo ma non giacenti sulla stessa retta di azione.
	$\Rightarrow$ costituistono una coppia di forze.\\\\ 
	Il \textbf{bracco} di ciascuna forza è: $\frac{\overline{AB}}{2} \sin \vartheta$\\\\
	$\Rightarrow$ il momento torcente $\tau$ è dato da:
	$$\tau = \underbrace{F_{BC} \frac{\overline{AB}}{2} \sin \vartheta}_{\#_1} + \underbrace{F_{AD} \frac{\overline{AB}}{2} \sin \vartheta}_{\#_2} = I \; A \; B \; \sin \vartheta$$
	Dove $A$ indica l'\textbf{area della spira}.\\
	Se definiamo $\vec{A} = A \hat{n}$ e riscriviamo l'equazione soprastante otteniamo:
	$$\Rightarrow \vec{\tau} = I \vec{A} \times \vec{B}$$
	e chiamando $\vec{\mu} = I \vec{A}$ \textbf{Momento Di Dipolo Magnetico} otteniamo infine che:
	$$\boxed{\vec{\tau} = \vec{\mu} \times \vec{B}}$$
	La relazione $\vec{\tau} = \vec{\mu} \times \vec{B}$ \textbf{vale per spire di forma qualsiasi}.\\\\
	Ogni spira percorsa da corrente è caratterizzata da un vettore $\vec{\mu} = I \vec{A}$ momento di dipolo magnetico.\\\\
	Il verso di rotazione tende ad allineare $\vec{\mu}$ con il campo $\vec{B}$
	\begin{figure}[h]
		
		\begin{subfigure}{0.55\textwidth}
			\includegraphics[scale=0.55]{117.png} 
			\caption{}
			\label{fig:im-117}
		\end{subfigure}
		\begin{subfigure}{0.55\textwidth}
			\includegraphics[scale=0.55]{118.png}
			\caption{}
			\label{fig:im-118}
		\end{subfigure}
		\begin{subfigure}{0.55\textwidth}
			\includegraphics[scale=0.55]{119.png}
			\caption{}
			\label{fig:im-119}
		\end{subfigure}
		
		\caption{}
		\label{fig:image13}
	\end{figure}
	Siamo in analogia con il comportamento di un ago magnetico ossia di un piccolo dipolo magnetico:
	\begin{figure}[h]
		
		\begin{subfigure}{0.55\textwidth}
			\includegraphics[scale=0.55]{120.png} 
			\caption{}
			\label{fig:im-120}
		\end{subfigure}
		\begin{subfigure}{0.55\textwidth}
			\includegraphics[scale=0.55]{121.png}
			\caption{}
			\label{fig:im-121}
		\end{subfigure}
		\begin{subfigure}{0.55\textwidth}
			\includegraphics[scale=0.55]{122.png}
			\caption{}
			\label{fig:im-122}
		\end{subfigure}
		
		\caption{Analogia con ago magnetico}
		\label{fig:image14}
	\end{figure}
	Agli effetti magnetici una spira piana di area $A$ e percorsa da una corrente $I$ equivale ad un \textbf{dipolo magnetico} di \textbf{momento magnetico di dipolo} $\vec{\mu} = I \, \vec{A}$
	\begin{center}
		\begin{figure}[h]
			\centering
			\includegraphics[scale=0.45]{123.png}
			\caption{}
			\label{fig:im-123}
		\end{figure}
	\end{center}
	È immediata l’analogia con il caso elettrico ricordando il momento di dipolo elettrico
	\begin{center}
		\begin{figure}[h]
			\centering
			\includegraphics[scale=0.45]{124.png}
			\caption{}
			\label{fig:im-124}
		\end{figure}
	\end{center}
	$$\vec{\tau} = \vec{P} \times \vec{E}$$
	Sempre per l’analogia formale:
	$$U = -\vec{P} \cdot \vec{E} \qquad \qquad \vec{\mathbf{U}} = -\vec{\mu} \cdot \vec{B}$$
	\textbf{Concetti già sviluppati per i dipoli elettrici potranno essere applicati a situazioni magnetiche}.
	\subsection{Moti della Particella}
	\subsection{Campo Magnetico e Fili Conduttori}
	\subsection{Momento Magnetico}
	\subsection{Legge di Biot-Savat e le sue Applicazioni}
	\subsection{Legge di Ampere}
	\subsection{Legge di Ampere-Maxwell}
	\subsection{Flusso Magnetico}
\end{document}